% =========================================================
% Density of the Rational Numbers
% =========================================================

\section{Density of the Rational Numbers}
\label{sec:density-q}

\subsection{Definitions and Theorems}

Two different ideas are often confused when one first studies the rational numbers.

\begin{itemize}
  \item The \emph{Archimedean property} says that the natural numbers grow beyond every fixed
  number.
  \item The \emph{density of $\mathbb{Q}$} says that between any two distinct real numbers there is a
  rational number.
\end{itemize}

The first controls large-scale size; the second controls local refinement of intervals.

\begin{theorem}[Archimedean property of $\mathbb{Q}$]
\label{thm:archimedean-q}
For every $x\in\mathbb{Q}$ there exists $n\in\mathbb{N}$ such that
\[
n>x.
\]
Equivalently, for every $x>0$ there exists $n\in\mathbb{N}$ such that
\[
\frac{1}{n}<x.
\]
\end{theorem}

\begin{remark}
This theorem says that the reciprocals of natural numbers become arbitrarily small.  It rules out the
existence of positive infinitesimals in $\mathbb{Q}$.
\end{remark}

\begin{theorem}[Density of $\mathbb{Q}$ in $\mathbb{R}$]
\label{thm:q-dense-in-r}
If $x,y\in\mathbb{R}$ and $x<y$, then there exists $q\in\mathbb{Q}$ such that
\[
x<q<y.
\]
\end{theorem}

\begin{proof}
Choose $n\in\mathbb{N}$ so large that
\[
\frac{1}{n}<y-x.
\]
By the Archimedean property there exists an integer $m$ such that
\[
m>nx.
\]
Choose the least such integer $m$.  Then $m-1\le nx$, and hence
\[
x<\frac{m}{n}\le x+\frac{1}{n}<y.
\]
Therefore $q=m/n$ is rational and satisfies $x<q<y$.
\end{proof}

\begin{corollary}
\label{cor:infinitely-many-rationals-between}
Between any two distinct real numbers there are infinitely many rational numbers.
\end{corollary}

\begin{proof}
Once one rational number is found between $x$ and $y$, the theorem can be applied again to the
subintervals that remain.  Repeating the argument indefinitely yields infinitely many rationals in the
interval.
\end{proof}

\subsection{Consequences}

\begin{remark}
Density is one of the reasons the rational numbers are so useful in analysis.  Even when a number is
not rational, it can still be approximated arbitrarily well by rational numbers.
\end{remark}
